\subsection{Characters of a commutative Banach algebra}

THEOREM 1. --- Let A be a Banach algebra and let \(\chi\colon\mathbf{A}\to\mathbf{C}\) be an algebra homomorphism (cf.~I, p.~9). Then \(\chi\) is continuous, with norm at most 1. If A is unital and if \(\chi\) is a unital homomorphism, then \(\chi\) has norm 1.

Let us prove that \(|\chi(x)| \leqslant \|x\|\) for all \(x \in \mathbf{A}\) . Replacing A, if necessary, by the Banach algebra generated by \(x\) , we may assume that A is commutative; then \(\chi \in \mathbf{X}'(\mathbf{A})\) . For every \(x \in \mathbf{A}\) , one has \(\chi(x) \in \mathrm{Sp}_{\mathrm{A}}'(x)\) (I, p.~9, no.~7), hence \(|\chi(x)| \leqslant \varrho(x) \leqslant \|x\|\) (I, p.~28, cor. 5), whence the first assertion.

If, moreover, A is unital, the equality \(\chi(1)=1\) implies that \(\|\chi\|\geqslant1\) , whence the desired equality.

Remark. --- Let A be a commutative Banach algebra. It follows from this theorem that the topology of pointwise convergence on \(X'(A)\) coincides with the restriction to \(X'(A)\) of the weak topology \(\sigma(A', A)\) of the dual \(A'\) of A.

COROLLARY. --- Let A be a commutative Banach algebra. The space X'(A) is compact. The space X(A) is locally compact, and is compact if A has a unit element.

Let A' be the dual of the Banach space A. Its unit ball \(A_{1}^{\prime}\) is weakly compact (EVT, III, p.~17, cor. 2). By th. 1, we have \(\mathsf{X}^{\prime}(\mathsf{A}) \subset \mathsf{A}_{1}^{\prime}\) . Moreover, \(\mathsf{X}^{\prime}(\mathsf{A})\) is closed in A' for the weak topology, since it is the intersection of the weakly closed sets

\[
\mathrm{X} _ {x, y} = \left\{f \in \mathrm{A} ^ {\prime} \mid f (x y) - f (x) f (y) = 0 \right\} \quad (\text { for } x, y \in \mathrm{A}).
\]

It follows that \(X'(A)\) is compact and that \(X(A) = X'(A) - \{0\}\) is locally compact.

If A has a unit element 1, then \(\mathsf{X}(\mathsf{A})\) is the set of \(\chi \in \mathsf{X}'(\mathsf{A})\) such that \(\chi(1) = 1\) , and is therefore a closed subset, and hence compact, of \(\mathsf{X}'(\mathsf{A})\) .

Every compact space is homeomorphic to X(A) for a suitable commutative unital Banach algebra A (cf.~I, p.~32, cor. 2).

THEOREM 2. --- Let A be a commutative Banach algebra. The map \(\chi \mapsto \operatorname{Ker}(\chi)\) is a bijection from \(\mathsf{X}(\mathsf{A})\) onto the set \(\mathrm{J}(\mathrm{A})\) of the regular maximal ideals of A.

Let I be a regular maximal ideal of A. It is closed (I, p.~22, cor. 2), hence A/I is a Banach algebra. Since it is a field (A, VIII, p.~426, prop. 2), the Gelfand--Mazur theorem (I, p.~26, cor. 2) implies that A/I has dimension 1 over C. Hence I has codimension 1 in A. The theorem then follows from lemma 3 of I, p.~10.

It may happen that a nonzero commutative Banach algebra A, without unit element, has no maximal ideal; then \(X'(A)\) is reduced to \(\{0\}\) (cf.~I, p.~186, exercise 31).

This theorem shows that the sets \(J(A)\) and \(\hat{A}\) from I, p.~11, which are identified since A is commutative, can also be identified with the set \(X(A)\) . There is therefore reason to consider on \(X(A)\) the weak topology and the Jacobson topology, which is coarser (I, p.~15, prop. 5). The closed sets \(V(M)\) of the Jacobson topology, for \(M \subset A\) , are the sets

\[
\{\chi \in \mathrm{X} (\mathrm{A}) \mid \chi (\mathrm{M}) = 0 \}
\]

of \(\mathsf{X}(\mathsf{A})\) , which will sometimes still be denoted \(\mathsf{V}(\mathsf{M})\) . Similarly, for any subset \(\mathsf{M}\) of \(\mathsf{X}(\mathsf{A})\) , we shall denote by \(\Upsilon(\mathsf{M})\) the ideal intersection of the kernels of the \(\chi \in \mathsf{M}\) ; it is equal to the set \(\Upsilon(\mathsf{M}')\) defined in I, p.~13 for the subset \(\mathsf{M}' \subset \mathsf{J}(\mathsf{A})\) corresponding to M when one identifies \(\mathsf{J}(\mathsf{A})\) and \(\mathsf{X}(\mathsf{A})\) .

When a topological notion is used in X(A) without specifying which topology is meant, it will always be the weak topology.

